\documentclass[11pt,a4paper]{article}
\usepackage{DDTA_DERMATRIAGE_GUIDED_REVIEW_STYLE_R1}
\providecommand{\tightlist}{\setlength{\itemsep}{0pt}\setlength{\parskip}{0pt}}
\begin{document}
\selectlanguage{italian}
\begin{titlepage}
\thispagestyle{empty}
\vspace*{10mm}
{\sffamily\bfseries\color{DDTANavy}\fontsize{15}{18}\selectfont Documentation-Driven Threat Analysis}\par
\vspace{5mm}
{\sffamily\bfseries\color{DDTANavy}\fontsize{27}{31}\selectfont DDTA Documentation Authoring Guide}\par
\vspace{2mm}
{\sffamily\color{DDTABlue}\fontsize{17}{21}\selectfont R5 Working Candidate - A5 Holdout Synthesis}\par
\vspace{10mm}
\begin{tcolorbox}[enhanced,arc=3pt,boxrule=0pt,colback=DDTANavy,left=12pt,right=12pt,top=10pt,bottom=10pt]
{\color{white}\sffamily\bfseries Current research artifact}\par
{\color{white!88}\small Repository baseline: \texttt{0e60754d21aa24ea487f3f60803b6b0cce8d2e2b}. Methodology authority remains R4 unless this document explicitly states otherwise.}
\end{tcolorbox}
\vfill
{\sffamily\small\color{DDTAGray} 5 settembre 2026}
\end{titlepage}
\tableofcontents
\clearpage
\section{DDTA Documentation Authoring Guide - R5 Working
Candidate}\label{ddta-documentation-authoring-guide---r5-working-candidate}

\textbf{Status:} WORKING CANDIDATE / A5-HOLDOUT SYNTHESIS /
NON-AUTHORITATIVE\\
\textbf{Methodology authority retained:}
\texttt{methodology/DDTA\_DOCUMENTATION\_BA\_AUTHORING\_GUIDE\_R4.tex}\\
\textbf{Repository baseline used for this synthesis:}
\texttt{0e60754d21aa24ea487f3f60803b6b0cce8d2e2b}\\
\textbf{Holdout evidence:} DermaTriage documentation review through A5
cumulative closure\\
\textbf{Date:} 5 September 2026

\begin{quote}
This revision is a complete forward rewrite of the
documentation-authoring guidance after the DermaTriage A5 holdout work.
It is deliberately \textbf{not promoted} to methodology authority by
this artifact. R4 remains authoritative until this candidate is
reviewed, regression-tested and explicitly promoted.
\end{quote}

\subsection{1. Purpose and thesis
boundary}\label{1-purpose-and-thesis-boundary}

DDTA documents the \textbf{minimum sufficient governed meaning} of a
project so that people can understand what the project is responsible
for, what material commitments it makes, what operational obligations
follow from those commitments, and what properties must hold, while
preserving explicit gaps instead of inventing missing detail.

The thesis scope is intentionally narrower than the full universe of
requirements engineering. The documentation metamodel keeps
\texttt{SpecializedRequirement} abstract because governed
specializations may exist for quality, performance, safety, reliability
and other concerns. The concrete specialization developed and evaluated
in depth by the thesis is \texttt{SecurityRequirement}, because the
research objective is security analysis over governed documentation.

Accordingly:

\begin{verbatim}
Requirement [abstract]
|-- FunctionalRequirement
`-- SpecializedRequirement [abstract]
    `-- SecurityRequirement       <- concrete specialization in thesis scope
\end{verbatim}

A non-security specialization discovered in a real project is
\textbf{not discarded}. Its governed meaning is preserved and its
existence is evidence that the abstract \texttt{SpecializedRequirement}
layer is useful. Designing a complete concrete subtype family for that
concern is outside the thesis scope unless a later explicit research
decision reopens it.

\subsection{2. Epistemic discipline: what is allowed to become project
truth}\label{2-epistemic-discipline-what-is-allowed-to-become-project-truth}

\subsubsection{2.1 Authority}\label{21-authority}

\textbf{Authority} answers: \emph{which artifact is allowed to define
project meaning at this moment?}

Before authoring or changing any DDTA element:

\begin{enumerate}
\def\labelenumi{\arabic{enumi}.}
\tightlist
\item
  identify the current governed baseline;
\item
  identify the authorized source set;
\item
  distinguish current, candidate, working, historical and superseded
  artifacts;
\item
  preserve source/revision provenance;
\item
  do not use recency as a substitute for authority;
\item
  do not allow Base Analysis, tests, code, tooling or threat-analysis
  output to become project authority automatically.
\end{enumerate}

Normal direction:

\begin{verbatim}
governed project documentation
        -> Base Analysis
        -> security/threat analysis
        -> finding / clarification / change candidate
        -> governed review
        -> updated project documentation
        -> rebuilt Base Analysis
\end{verbatim}

\subsubsection{2.2 Governed meaning}\label{22-governed-meaning}

A proposition is \textbf{governed} when the project has enough authority
and semantic commitment for that proposition to be treated as part of
the project contract.

A source may accurately describe a current implementation without
governing that implementation choice for the future. Therefore:

\begin{verbatim}
source-observed fact != automatically governed requirement
\end{verbatim}

Example from DermaTriage:

\begin{verbatim}
RETRAINING_THRESHOLD=50
\end{verbatim}

is useful realization/configuration evidence. The governed meaning
identified during the holdout is that classifier adaptation is
accumulation-gated by a threshold. The exact value \texttt{50} is a
current documented binding whose fixed normative status and change
authority are not specified.

\subsubsection{2.3 Semantic owner}\label{23-semantic-owner}

The \textbf{semantic owner} is the project
responsibility/decision/requisite to which a meaning belongs.

Before creating a new element ask:

\begin{verbatim}
Who owns this meaning?
\end{verbatim}

Do not infer ownership merely because two facts appear in the same
component, API response, table or sentence.

\subsubsection{2.4 NOT SPECIFIED}\label{24-not-specified}

\texttt{NOT\ SPECIFIED} is a valid documented result when the source
does not establish a material distinction sufficiently.

It is preferable to:

\begin{verbatim}
NOT SPECIFIED
\end{verbatim}

than to create a precise but invented rule.

\subsubsection{2.5 STOP}\label{25-stop}

\texttt{STOP} means that no additional independently governed meaning
has been justified at the current layer and scope. It does \textbf{not}
mean the project has no more implementation detail.

Examples:

\begin{itemize}
\tightlist
\item
  DermaTriage MR-02 stops at MR because specialist destination semantics
  are not sufficiently governed downstream.
\item
  A FunctionalRequirement may stop at FR when no independently governed
  specialized property is established.
\end{itemize}

\subsection{3. L1-L4 separation}\label{3-l1-l4-separation}

DDTA separates four layers:

\begin{longtable}[]{@{}lll@{}}
\toprule\noalign{}
Layer & Meaning & Typical contents \\
\midrule\noalign{}
\endhead
\bottomrule\noalign{}
\endlastfoot
L1 & conceptual metamodel & MR, Decision, Requirement, FR, SR,
SecurityRequirement, cardinalities and semantic invariants \\
L2 & authoring/representation contract & headings, fields, controlled
terminology, normative-clause representation, registries \\
L3 & realization/model principles & schemas, validators, model loaders,
derived indexes \\
L4 & tool support/conformance & editor assistance, checks, projections,
reports \\
\end{longtable}

A convenient YAML key or Markdown heading is not automatically a new L1
concept.

\subsection{4. Normal DDTA authoring versus holdout
reconstruction}\label{4-normal-ddta-authoring-versus-holdout-reconstruction}

\subsubsection{4.1 Normal authoring}\label{41-normal-authoring}

The normal process starts from project responsibility and governance:

\begin{verbatim}
problem / responsibility
        -> authority
        -> MacroRequirement
        -> Decision
        -> FunctionalRequirement
        -> specialization review
        -> SecurityRequirement where governed
\end{verbatim}

\subsubsection{4.2 Existing-documentation
validation}\label{42-existing-documentation-validation}

Reconstruction from traditional documentation is a \textbf{validation
protocol}, not the normal production method:

\begin{verbatim}
existing project documentation
        -> recover only supported project meaning
        -> apply normal DDTA authoring gates
        -> candidate governed documentation
        -> compare preserved / clarified / unresolved / lost
\end{verbatim}

During a holdout, Base Analysis must not be consulted to decide what
project meaning to introduce. A holdout is valuable precisely because it
may expose a representability limitation.

\subsection{5. End-to-end authoring
sequence}\label{5-end-to-end-authoring-sequence}

The current working sequence is:

\begin{verbatim}
A0 Authority/source gate
A1 Project problem framing
A2 MacroRequirement / D1
A3 Semantic sufficiency
A4 Decision / D2
A5 FunctionalRequirement / D3
A6 Requirement split
PRE-A7 scope + documentation<->BA consistency check
A7 SpecializedRequirement review
A8 SecurityRequirement review
A9 Cross-MR / consumed-service review
A10 Terminology / information / parameter bindings
A11 Downstream propagation / semantic regression
A12 Completeness / promotion readiness
ONLY THEN accepted Base Analysis
THEN security/threat analysis such as STRIDE / STRIDE-AI
THEN governed feedback into documentation
\end{verbatim}

Local coherent-unit/split checks may occur before A6, but A6 is the
project-wide regression of those local decisions.

\subsection{6. A0 - Authority and source
gate}\label{6-a0---authority-and-source-gate}

Before deriving meaning record:

\begin{itemize}
\tightlist
\item
  source package identity and hash where available;
\item
  principal semantic source;
\item
  evidence/realization sources;
\item
  excluded or non-authoritative reconstructions;
\item
  current repository/methodology baseline.
\end{itemize}

For a holdout, distinguish at least:

\begin{verbatim}
project authority
principal semantic source
evidence / realization source
historical reconstruction
analysis artifact
\end{verbatim}

DermaTriage example:

\begin{itemize}
\tightlist
\item
  authoritative package:
  \texttt{DermaTriage-Docs-20260830T152637Z-1-001.zip};
\item
  package SHA-256:
  \texttt{E9ED2C507BEFB95F54A52084687CD1E8798863AE81CF69D09568864D8CBF280E};
\item
  OR2 Architecture Document: principal semantic source;
\item
  test/training/dataset/configuration documents: supporting evidence
  unless they clearly govern project meaning;
\item
  previous DDTA/BA reconstructions: not project authority.
\end{itemize}

\subsection{7. A1 - Project problem
framing}\label{7-a1---project-problem-framing}

Start from the problem responsibility, not the implementation stack.

Bad framing:

\begin{verbatim}
DermaTriage uses EfficientNet, Qwen, ChromaDB and BioMistral.
\end{verbatim}

Better framing:

\begin{verbatim}
DermaTriage supports early triage of dermatological cases concerning potentially oncological skin lesions,
using available case information to determine a priority of care and support routing toward appropriate specialist assistance.
\end{verbatim}

The framing may expose unresolved boundaries. In DermaTriage, authority
for definitive clinical diagnosis remains outside the established
governed meaning.

\subsection{8. A2 - MacroRequirement}\label{8-a2---macrorequirement}

A \texttt{MacroRequirement} governs one stable macro responsibility.

Recommended fields:

\begin{verbatim}
Title
Intent
Context
Stakeholders
Scope
Assumptions / Constraints   [only when branch-wide]
dependsOn                  [only when genuinely governed]
\end{verbatim}

\subsubsection{8.1 MR split test}\label{81-mr-split-test}

Split when two responsibilities can exist, evolve or be governed
independently.

DermaTriage produced four macro responsibilities:

\begin{itemize}
\tightlist
\item
  MR-01 triage evaluation;
\item
  MR-02 specialist routing support;
\item
  MR-03 clinical review management;
\item
  MR-04 controlled adaptation from accumulated clinical review evidence.
\end{itemize}

Clinical review and future adaptation were split because review can
exist without retraining, and adaptation policy can change while review
capture remains stable.

\subsubsection{8.2 MR STOP}\label{82-mr-stop}

If the macro responsibility is governed but no material downstream
choice/obligation is sufficiently specified, STOP at MR.

DermaTriage MR-02 is the canonical example: an indication of specialist
destination is present, but complete selection/routing semantics are not
sufficiently defined.

\subsection{9. A3 - Semantic
sufficiency}\label{9-a3---semantic-sufficiency}

Ask for the \textbf{smallest critical proposition} whose ambiguity would
create materially different project meanings.

Possible dispositions:

\begin{verbatim}
AFFIRMED
DENIED
NOT SPECIFIED
CONFLICTING
REWORK
\end{verbatim}

Numerical precision does not imply semantic sufficiency. A table may say
\texttt{24h}, \texttt{48h}, \texttt{72h}, \texttt{7\ days}; if the
documentation does not say what event must occur within those times,
from which trigger the interval runs, and who owns compliance, the
timing semantics remain \texttt{NOT\ SPECIFIED}.

\subsection{10. A4 - Decision}\label{10-a4---decision}

A \texttt{Decision} governs one material commitment that narrows exactly
one MR.

Recommended structure:

\begin{verbatim}
Context
Decision
Consequences
\end{verbatim}

A Decision should survive reasonable implementation substitution.

\subsubsection{10.1 Decision kinds}\label{101-decision-kinds}

Use the existing kinds:

\begin{itemize}
\tightlist
\item
  \texttt{SELECTABLE} - materially different alternatives exist;
\item
  \texttt{NECESSITY-CONSTRAINED} - project constraints force the
  commitment;
\item
  \texttt{DEFAULT} - a default position is governed while change remains
  possible.
\end{itemize}

Do not create fake Decision wrappers merely because the hierarchy
expects a Decision before FR.

\subsubsection{10.2 Decision versus operational
rule}\label{102-decision-versus-operational-rule}

A Decision answers \emph{which material project position is adopted}. A
FunctionalRequirement answers \emph{what operational behavior is
required because of that position}.

DermaTriage example:

\begin{verbatim}
DEC-02: adopt P1-P4 as operational priority scale
FR-02: derive the concrete P-scale level from urgency/confidence
\end{verbatim}

\subsubsection{10.3 Numeric neutralization at Decision
level}\label{103-numeric-neutralization-at-decision-level}

When a candidate Decision is dominated by a concrete number, replace the
number with a symbolic value and ask whether the policy still exists.

Example:

\begin{verbatim}
Every 10 corrections -> evolve prompt
\end{verbatim}

Neutralize:

\begin{verbatim}
When accumulated relevant evidence reaches threshold N -> evolve prompt
\end{verbatim}

The surviving Decision is the accumulation-gated adaptation strategy.
\texttt{10} is not what makes the Decision exist.

\subsection{11. A5 -
FunctionalRequirement}\label{11-a5---functionalrequirement}

A \texttt{FunctionalRequirement} is one independently assessable
operational obligation under exactly one Decision.

A strong FR states enough to determine:

\begin{itemize}
\tightlist
\item
  triggering/applicability condition;
\item
  responsible project behavior;
\item
  input/output/state meaning needed to assess it;
\item
  material context binding;
\item
  failure or non-satisfaction semantics where governed;
\item
  deliberate exclusions.
\end{itemize}

\subsubsection{11.1 One FR does not equal one
sentence}\label{111-one-fr-does-not-equal-one-sentence}

\begin{verbatim}
1 Requirement != 1 textual sentence
1 Requirement  = 1 coherent normative obligation
\end{verbatim}

A mapping table may contain several normative clauses and still be one
FR if all rows jointly express one coherent transformation.

DermaTriage FR-08:

\begin{verbatim}
P1 -> HIGH
P2 -> HIGH
P3 -> MEDIUM
P4 -> LOW
\end{verbatim}

is one FR, not four requirements.

\subsubsection{11.2 Do not merely restate the
Decision}\label{112-do-not-merely-restate-the-decision}

A Decision such as ``use clinical disagreement as qualifying evidence''
may still need an FR that makes the operational qualification rule
assessable.

\subsubsection{11.3 Responsibility boundary before
FR}\label{113-responsibility-boundary-before-fr}

Do not assign to the project a behavior owned by a consumed service, a
clinician or another responsibility.

DermaTriage FR-03 records/correlates a clinical review; it does not
claim that DermaTriage performs the clinical judgment.

\subsubsection{11.4 Contextual binding}\label{114-contextual-binding}

When a relation must remain tied to the same case/request/outcome to
avoid a materially different result, preserve that binding in FR
semantics.

\subsection{12. Decision-to-Requirement completeness
counterexample}\label{12-decision-to-requirement-completeness-counterexample}

Before declaring STOP below a Decision, ask:

\begin{verbatim}
Can all current downstream Requirements be satisfied while this Decision is still violated?
\end{verbatim}

\begin{itemize}
\tightlist
\item
  If \textbf{NO}, STOP may be justified.
\item
  If \textbf{YES}, downstream coverage is incomplete; search for the
  missing FR/SR rather than reopening the Decision automatically.
\end{itemize}

DermaTriage DEC-05 exposed this rule. FR-04 and FR-05 could both be
satisfied while one adaptation path incorrectly triggered the other.
FR-11 was therefore required to govern lifecycle independence.

\subsection{13. Parameter Governance
Boundary}\label{13-parameter-governance-boundary}

Do not create a requirement or document for every number.

The correct sequence is:

\begin{verbatim}
source value
    -> semantic owner
    -> lifecycle / purpose
    -> requirement/property classification
    -> neutralize concrete value to symbolic parameter
    -> semantic materiality test
    -> exact-value governance test
    -> governed semantic parameter OR realization/configuration STOP
\end{verbatim}

\subsubsection{13.1 Gate Q1 - does meaning survive replacement by
N?}\label{131-gate-q1---does-meaning-survive-replacement-by-n}

If replacing the number by \texttt{N} destroys the governed proposition,
the number may be part of the semantic rule itself rather than a
replaceable parameter.

Example: the transformation \texttt{P2\ -\textgreater{}\ HIGH} cannot be
neutralized to \texttt{P2\ -\textgreater{}\ X} without losing the
governed meaning.

\subsubsection{13.2 Gate Q2 - can parameter variation materially change
governed
behavior?}\label{132-gate-q2---can-parameter-variation-materially-change-governed-behavior}

If no, the value is ordinary technical configuration and should normally
stop below the semantic model.

If yes, preserve a semantic parameter.

\subsubsection{13.3 Gate Q3 - is the concrete value itself
governed?}\label{133-gate-q3---is-the-concrete-value-itself-governed}

\begin{itemize}
\tightlist
\item
  YES: preserve the exact value normatively.
\item
  NO / NOT SPECIFIED: preserve the symbolic parameter and record the
  concrete value as a current binding.
\end{itemize}

DermaTriage examples:

\begin{longtable}[]{@{}lrl@{}}
\toprule\noalign{}
Semantic parameter & Current documented binding & Exact value
governance \\
\midrule\noalign{}
\endhead
\bottomrule\noalign{}
\endlastfoot
\texttt{PromptEvolutionThreshold=N} & 10 & NOT SPECIFIED \\
\texttt{ClassifierAdaptationThreshold=N} & 50 & NOT SPECIFIED \\
\texttt{PromptEvidenceWindowSize=N} & 20 & NOT SPECIFIED \\
\texttt{AccuracyDegradationTolerance=T} & 5\% & quantitative binding
incomplete \\
\texttt{RollbackAccuracyDegradationThreshold=R} & 5\% & quantitative
binding incomplete \\
\end{longtable}

The two \texttt{5\%} values are \textbf{not} one parameter merely
because the literal is equal. Semantic owner, lifecycle and purpose
differ.

\subsection{14. Information and realization bridge
taxonomy}\label{14-information-and-realization-bridge-taxonomy}

DDTA must preserve enough traceability to connect governed semantics to
implementation without making implementation identifiers normative by
default.

Observed bridge shapes:

\begin{verbatim}
1. semantic concept
   -> data/state encoding

2. semantic parameter
   -> configuration value

3. semantic reference
   -> interface identifier

4. semantic reference
   -> governed transformation
   -> semantic output
   -> realization encoding
\end{verbatim}

DermaTriage examples:

\begin{verbatim}
ClinicianDisagreement
    -> source-observed encoding: agrees == False

ClassifierAdaptationThreshold=N
    -> current binding: 50

ClinicallyCorrectedPriority
    -> P1/P2/P3/P4 mapping
    -> ClassifierUrgencyTarget
    -> enum/string/class-index realization
\end{verbatim}

The semantic concept remains stable even if a field name or enum
representation changes.

\subsection{15. A6 - Requirement split}\label{15-a6---requirement-split}

A6 is the project-wide coherent-unit regression.

Split a Requirement when a part can be introduced, revised, retired,
assessed, fail or be revalidated independently without changing the
normative identity of the remaining parts.

Strong split indicators:

\begin{itemize}
\tightlist
\item
  different semantic owner;
\item
  materially distinct failure mode;
\item
  independent lifecycle/review;
\item
  different specialization/security scope;
\item
  independent change/revalidation.
\end{itemize}

Do not split merely because there are multiple verbs, multiple mapping
rows or multiple test cases.

\subsection{16. Pre-A7 scope and Documentation\textless-\textgreater BA
consistency
check}\label{16-pre-a7-scope-and-documentation-ba-consistency-check}

Before A7 in the current thesis workflow perform a controlled
consistency check:

\begin{enumerate}
\def\labelenumi{\arabic{enumi}.}
\tightlist
\item
  verify that the documentation metamodel can preserve each discovered
  specialization meaning;
\item
  verify that the BA handoff can consume required project meaning
  without dictating document types;
\item
  preserve the thesis boundary: security specialization is the concrete
  research scope;
\item
  if a non-security specialization is discovered, preserve it without
  opening a full new subtype research program;
\item
  a pre-A7 BA probe, if used, is non-authoritative and exists only to
  test representability, not to decide project meaning.
\end{enumerate}

\subsection{17. A7 -
SpecializedRequirement}\label{17-a7---specializedrequirement}

A \texttt{SpecializedRequirement} is justified when an additional
concern-specific property strengthens exactly one parent FR and can be
governed independently from the ordinary functional behavior.

Semantic invariants:

\begin{itemize}
\tightlist
\item
  exactly one parent FR;
\item
  conjunctive strengthening of the parent;
\item
  removing the specialization leaves the ordinary function meaningful;
\item
  no duplication of ordinary functional correctness;
\item
  subordinate positive work may be required;
\item
  an autonomous new capability belongs in a separate FR;
\item
  realization mechanism is not the property;
\item
  the specialization depends on the parent behavior.
\end{itemize}

\subsubsection{17.1 Scope-limited specialization
rule}\label{171-scope-limited-specialization-rule}

In this thesis:

\begin{verbatim}
specialization candidate
    -> security property? YES -> A8 SecurityRequirement
    -> security property? NO
         -> preserve governed property
         -> record extensibility evidence
         -> mark concrete non-security subtype OUTSIDE THESIS SCOPE
\end{verbatim}

DermaTriage DEC-07 is the main holdout evidence: sensitivity
non-degradation, false-low non-degradation and bounded accuracy
degradation are plausible quality specializations of FR-09, but they are
not security requirements and therefore do not justify expanding the
thesis into a general quality-requirements taxonomy.

\subsection{18. A8 -
SecurityRequirement}\label{18-a8---securityrequirement}

A \texttt{SecurityRequirement} is a concrete specialized requirement in
thesis scope.

Ask:

\begin{enumerate}
\def\labelenumi{\arabic{enumi}.}
\tightlist
\item
  is the security property governed by project authority?
\item
  is its scope tied to a precise governed behavior/meaning?
\item
  can it evolve/fail independently from other properties?
\item
  are we stating a property rather than merely a mechanism?
\item
  is the security failure mode identifiable in the normative clauses?
\end{enumerate}

Never infer:

\begin{verbatim}
Confidentiality -> TLS
Integrity -> digital signature
Authorized provenance -> certificate
Authentication -> X-API-Key
\end{verbatim}

unless the mechanism itself is a governed commitment.

Existing project documentation may already govern SecurityRequirements
before threat analysis. Threat analysis may later produce
\textbf{candidate} security obligations, but findings do not become
SecurityRequirements automatically.

\subsection{19. A9 - Cross-MR and consumed-service
semantics}\label{19-a9---cross-mr-and-consumed-service-semantics}

Distinguish:

\begin{verbatim}
ownership
consumption
reference
dependency
second parent
\end{verbatim}

They are not synonyms.

A FunctionalRequirement may reference a concept governed by another
Decision/MR without acquiring a second parent.

DermaTriage DEC-11 consumes the P-scale domain governed under MR-01
while remaining owned by MR-04.

Consumed services do not become project-owned merely because their
outputs are used.

\subsection{20. A10 - Terminology, information, parameter and interface
bindings}\label{20-a10---terminology-information-parameter-and-interface-bindings}

A10 consolidates the bridge ledger accumulated during authoring.

For each important datum classify it as one of:

\begin{enumerate}
\def\labelenumi{\arabic{enumi}.}
\tightlist
\item
  governed normative meaning;
\item
  source-observed vocabulary/representation;
\item
  unresolved binding / NOT SPECIFIED;
\item
  current realization/configuration binding.
\end{enumerate}

A10 should maintain structured registries rather than one document per
value or field.

Do not invent requiredness, cardinality, type, missing-value policy,
reset policy or change authority merely to make a data model look
complete.

\subsection{21. A11 - Downstream propagation and semantic
regression}\label{21-a11---downstream-propagation-and-semantic-regression}

Recheck that upstream corrections are preserved through every descendant
and reference.

Questions:

\begin{itemize}
\tightlist
\item
  did any downstream artifact silently narrow or widen an upstream
  meaning?
\item
  did a current implementation value become a fixed normative
  requirement accidentally?
\item
  did a cross-MR reference become ownership?
\item
  did a gap disappear because of wording rather than evidence?
\item
  can a Decision still be violated while all descendants pass?
\item
  do canonical terms still denote the same semantic referent?
\end{itemize}

\subsection{22. A12 - Completeness and promotion
readiness}\label{22-a12---completeness-and-promotion-readiness}

A documentation baseline can be promoted for BA use when:

\begin{itemize}
\tightlist
\item
  material semantic owners are established;
\item
  required dependencies/cross-references are governed;
\item
  blocking gaps are resolved or explicitly excluded from the declared
  scope;
\item
  non-blocking NOT SPECIFIED meanings are preserved;
\item
  specialization/security review is complete for thesis scope;
\item
  A10 bindings are coherent;
\item
  A11 semantic regression passes;
\item
  governance explicitly authorizes the baseline as BA source.
\end{itemize}

Promotion is not ``everything is fully specified''. Promotion means the
baseline is \textbf{sufficiently governed for the declared scope}.

\subsection{23. Handoff to Base
Analysis}\label{23-handoff-to-base-analysis}

The handoff should include:

\begin{itemize}
\tightlist
\item
  authority/baseline identity;
\item
  pinned source revision/source set;
\item
  canonical governed subjects/referents;
\item
  material behaviors, relations, conditions and states;
\item
  specialized/security properties in scope;
\item
  parameter semantics and current bindings where relevant;
\item
  diagnostics, NOT SPECIFIED meanings and unresolved bindings;
\item
  change/revalidation context when material.
\end{itemize}

The BA then applies its own minimum representation criterion. It must
not retroactively dictate how the documentation should have been
written.

\subsection{24. STRIDE / STRIDE-AI and governed
feedback}\label{24-stride--stride-ai-and-governed-feedback}

The intended thesis lifecycle is:

\begin{verbatim}
GOVERNED DOCUMENTATION
        -> ACCEPTED BASE ANALYSIS
        -> STRIDE / STRIDE-AI
        -> FINDINGS
        -> candidate clarification / architecture change / accepted risk / SecurityRequirement
        -> GOVERNANCE REVIEW
        -> updated governed documentation if accepted
        -> rebuilt BA
\end{verbatim}

A threat finding is analytical evidence, not project truth.

\subsection{25. Extended worked example - DermaTriage adaptation
branch}\label{25-extended-worked-example---dermatriage-adaptation-branch}

\subsubsection{25.1 Source text}\label{251-source-text}

OR2 Architecture describes:

\begin{verbatim}
Prompt evolution trigger: every 10 doctor corrections
Prompt evidence method: sliding window of 20 most recent examples
Classifier retraining trigger: 50 doctor corrections where agrees == False
Label mapping: P1/P2 -> HIGH, P3 -> MEDIUM, P4 -> LOW
Deploy condition: sensitivity/false-low non-degradation and bounded accuracy degradation
Rollback: current implementation reports automatic rollback beyond 5% accuracy degradation
\end{verbatim}

\subsubsection{25.2 Recover MR meaning}\label{252-recover-mr-meaning}

The stable responsibility is not ``run PromptManager and EfficientNet''.
It is:

\begin{verbatim}
MR-04 - Controlled adaptation based on accumulated clinical review evidence
\end{verbatim}

\subsubsection{25.3 Recover Decisions}\label{253-recover-decisions}

The source supports distinct commitments:

\begin{itemize}
\tightlist
\item
  DEC-04 accumulation-gated adaptation;
\item
  DEC-05 separate adaptation paths by behavioral capability;
\item
  DEC-06 comparative pre-adoption qualification;
\item
  DEC-07 asymmetric quality prioritization;
\item
  DEC-08 classification adaptation reversibility;
\item
  DEC-09 bounded recent evidence for prompt evolution;
\item
  DEC-10 clinical disagreement as qualifying classifier-adaptation
  evidence;
\item
  DEC-11 corrected P-scale as classifier-supervision source.
\end{itemize}

\subsubsection{25.4 Recover FRs without hardcoding
implementation}\label{254-recover-frs-without-hardcoding-implementation}

\begin{verbatim}
FR-04 prompt adaptation activation at semantic threshold N
FR-05 classifier adaptation activation at semantic threshold N
FR-06 bounded recent prompt-evidence window
FR-07 clinical disagreement qualification
FR-08 corrected P-scale -> classifier target mapping
FR-09 comparative qualification
FR-10 rollback capability for materially degrading adopted adaptation
FR-11 path lifecycle independence
\end{verbatim}

\subsubsection{25.5 Preserve bindings rather than freezing
them}\label{255-preserve-bindings-rather-than-freezing-them}

\begin{verbatim}
PromptEvolutionThreshold=N           current N=10
ClassifierAdaptationThreshold=N      current N=50
PromptEvidenceWindowSize=N           current N=20
\end{verbatim}

The current literals are traceability evidence, not automatically
immutable normative constants.

\subsubsection{25.6 Preserve data/state encoding
separately}\label{256-preserve-datastate-encoding-separately}

\begin{verbatim}
ClinicianDisagreement
    governed semantic concept

agrees == False
    source-observed current encoding
\end{verbatim}

Changing the field name does not change the requirement if the same
semantic state remains represented.

\subsubsection{25.7 Preserve non-security specialization without
expanding thesis
scope}\label{257-preserve-non-security-specialization-without-expanding-thesis-scope}

DEC-07 points to quality properties that specialize FR-09. They are
preserved for A7 review. They are not reclassified as
SecurityRequirements and they do not force a new
\texttt{QualityRequirement} subtype in the current thesis.

\subsection{26. Practical authoring
worksheet}\label{26-practical-authoring-worksheet}

For every candidate element record:

\begin{longtable}[]{@{}ll@{}}
\toprule\noalign{}
Question & Result \\
\midrule\noalign{}
\endhead
\bottomrule\noalign{}
\endlastfoot
source evidence & exact source/section or governed predecessor \\
candidate meaning & concise project-semantic statement \\
semantic owner & MR / Decision / FR / specialization \\
material alternatives & what materially different reading/change is
excluded? \\
responsibility boundary & who actually owns the
behavior/judgment/property? \\
implementation-neutrality test & does the meaning survive technology
substitution? \\
split/coherent-unit test & can parts change or fail independently? \\
parameter/binding test & semantic parameter or current realization
value? \\
gap status & AFFIRMED / DENIED / NOT SPECIFIED / CONFLICTING \\
disposition & PASS / REWORK / STOP / ROUTE \\
methodology feedback & guide, metamodel, BA, tooling or project-source
pressure? \\
\end{longtable}

\subsection{27. Final authoring
checklist}\label{27-final-authoring-checklist}

Before promoting a documentation baseline ask:

\begin{enumerate}
\def\labelenumi{\arabic{enumi}.}
\tightlist
\item
  Is the source/authority boundary explicit?
\item
  Is every MR a stable responsibility rather than a technology bundle?
\item
  Does every Decision contain one material commitment?
\item
  Does every FR have exactly one honest parent Decision?
\item
  Can every FR be independently assessed without reconstructing hidden
  obligations from its parent?
\item
  Have mapping rows and normative clauses been kept in one Requirement
  when they form one coherent obligation?
\item
  Have independent obligations been split?
\item
  Have concrete numeric values passed semantic-owner and
  parameter-governance gates?
\item
  Are semantic concepts separated from current field/config/interface
  encodings?
\item
  Could all current descendants pass while a Decision is still violated?
\item
  Are cross-MR references distinguished from ownership and dependency?
\item
  Are non-security specialization candidates preserved without expanding
  the thesis scope unnecessarily?
\item
  Are SecurityRequirements properties rather than mechanisms?
\item
  Are unresolved meanings explicitly NOT SPECIFIED rather than silently
  invented?
\item
  Has project-wide propagation/regression been run?
\item
  Is the baseline explicitly authorized before Base Analysis begins?
\end{enumerate}

\subsection{28. Status of this R5 working
candidate}\label{28-status-of-this-r5-working-candidate}

This document incorporates holdout-supported refinements through
DermaTriage A5 closure, especially:

\begin{itemize}
\tightlist
\item
  semantic-owner-first numeric classification;
\item
  Parameter Governance Boundary;
\item
  information/parameter/reference/transformation bridge taxonomy;
\item
  Decision-to-Requirement counterexample test;
\item
  scope-limited handling of non-security specializations;
\item
  explicit documentation -\textgreater{} BA -\textgreater{} STRIDE
  -\textgreater{} governed feedback lifecycle.
\end{itemize}

It \textbf{does not} by itself promote these refinements into
methodology authority. Promotion requires review against R4, A6-A12
evidence, the pre-A7 documentation/BA consistency check, and regression
against the Facial Access corpus.

\end{document}
